\documentclass[10pt,a4paper]{amsart}
\usepackage[margin=19mm]{geometry}
\usepackage{amsmath,amssymb,mathtools}
\usepackage[scr=rsfs,cal=euler]{mathalfa}
\usepackage{xfrac}
\usepackage[unicode,hidelinks,bookmarks=false]{hyperref}
\newcommand{\ie}{i.e.\ }
\newcommand{\cf}{cf.\ }
\newcommand{\resp}{respectively\ }
\newcommand{\Subsection}{\subsection}
\providecommand{\nobreakdash}{\nobreak\hbox{-}}
\setlength{\emergencystretch}{2em}
\multlinegap0pt
\usepackage{amssymb}
\usepackage{mathtools}
\usepackage[shortlabels]{enumitem}
\usepackage{tikz}
\usepackage{xcolor}
\newtheorem{remark}{Remark}
\newtheorem{theorem}{Theorem}
\newtheorem{lemma}{Lemma}
\newcommand{\Cbracket}{B}            % bracket continuity constant
\newcommand{\Cmult}{C_{\cdot}}       % submultiplicativity constant
\newcommand{\Ctotal}{K}              % combined constant for BCH estimates
\newcommand{\Ctri}{C_{\triangle}}    % quasi-triangle constant
\newcommand{\pnorm}[1]{\left\vert\!\left\vert\!\left\vert #1 \right\vert\!\right\vert\!\right\vert} % equivalent p-norm
\title{Local Lie Theory in Quasi-Banach Lie Algebras: Convergence of the BCH Series and Geometric Implications}
\author{Nassim Athmouni, Mohsen Ben Abdallah, Mondher Damak, Marwa Ennaceur, Amel Jadlaoui, Lotfi Souden}
\date{Source: arXiv:2604.08265v1 (CC BY 4.0)}
\begin{document}
\maketitle

\noindent\emph{Source and licence.} Local Lie Theory in Quasi-Banach Lie Algebras: Convergence of the BCH Series and Geometric Implications. Version arXiv:2604.08265v1 (CC BY 4.0), distributed by arXiv under the Creative Commons Attribution 4.0 International licence, http://creativecommons.org/licenses/by/4.0/, as recorded in the arXiv licence metadata for this exact version and shown on the abstract page https://arxiv.org/abs/2604.08265v1. Full licence notice: \texttt{licenses/arxiv-2604.08265-CC-BY-4.0.txt}.\par\smallskip

\noindent\textit{Editorial scope.} Counted unit: the article's lemma lem:bch-lipschitz (source0.tex lines 609--624). The article's complete original proof of it is delivered with it (source0.tex lines 626--656, one \texttt{proof} environment of 31 lines). No mathematics is written, repaired or re-derived: the statement and the proof are the article's own lines, in the article's own notation, and no retained line is substituted or re-flowed.  Retained uncounted as the source-local context the proof needs: lines 160--171; lines 271--279; lines 440--470; lines 481--483; lines 556--560.  Nothing was omitted from any retained span, so the package carries no omission record.  No retained line contains a citation, so the wrapper carries no bibliography and no bibliography entry.  No retained line contains a figure environment, an image-inclusion command or drawing code, so no image is delivered and the package declares no asset dependency.  Editorial formatting only, in addition: an AMS \texttt{amsart} preamble that restates the article's own declarations and macros used by the retained lines, namely the environments \texttt{remark}, \texttt{theorem}, \texttt{lemma} on separate counters, together with the standard packages those lines need.  The retained environments print in this document as Remark 1, Theorem 1, Lemma 1 in order of appearance, whereas the article prints the same items with the numbering of its own full document.  The complete native source of the article as distributed in the arXiv e-print for this exact version is bundled unchanged as \texttt{sources/198147/source0.tex}.
\begin{remark}[Role of the constant $\Ctri$]\label{rem:role-ctri}
The Catalan-majorant argument alone gives the convergence radius $1/(4\Cbracket)$.  When
working with the $p$-norm equivalence, the factor $c_2 = 2^{1/p} = 2\Ctri$ is absorbed
into the convergence criterion as follows.  From \eqref{eq:pnorm-bound} one needs
$4\Cbracket r < 1$; the $c_2$-factor only appears in the \emph{value} of the sum
$\sum_n\pnorm{Z_n}^p$, not in the convergence condition.  Hence the combined constant
$\Ctotal = \Ctri\Cbracket$ represents a convenient but conservative upper bound: the true
radius is $1/(4\Cbracket)$, and the stated radius $1/(4\Ctotal)$ adds an extra safety
margin of $\Ctri$ arising from the norm equivalence.  We retain $\Ctotal$ throughout for
uniform presentation across the abstract and associative settings, but alert the reader
that the $\Ctri$-factor is not intrinsic to the BCH series itself.
\end{remark}

\begin{remark}[Bracket constant in associative embeddings]\label{rem:bracket-associative}
If $\mathfrak{g}$ is a Lie subalgebra of an associative quasi-Banach algebra
$(\mathcal{A},\|\cdot\|)$ with submultiplicativity constant $\Cmult$ and quasi-triangle
constant $\Ctri$, then the commutator bracket satisfies
\[
  \|[x,y]\| = \|xy - yx\| \le \Ctri(\|xy\| + \|yx\|) \le 2\Ctri\Cmult\|x\|\|y\|.
\]
Thus, in this setting, one may take $\Cbracket = 2\Ctri\Cmult$.
\end{remark}

\begin{theorem}[BCH convergence in quasi-Banach Lie algebras]\label{thm:BCH}
Let $(\mathfrak{g},[\cdot,\cdot],\|\cdot\|)$ be a quasi-Banach Lie algebra with
quasi-triangle constant $\Ctri$ and bracket continuity constant $\Cbracket$.  Then the
BCH series $Z(x,y) = \sum_{n\ge 1} Z_n(x,y)$ converges in the $d_p$-metric topology of
$\mathfrak{g}$, where $d_p(x,y) = \|x-y\|^p$ and $p = 1/\log_2(2\Ctri)$, for all $x,y$
satisfying
\begin{equation}\label{eq:BCH-radius}
  \|x\| + \|y\| < \frac{1}{4\Cbracket}.
\end{equation}
In particular, convergence holds on the symmetric domain $\|x\|,\|y\| < 1/(8\Cbracket)$.

A conservative uniform bound valid in all quasi-Banach settings is obtained by setting
$\Ctotal := \Ctri \cdot \Cbracket$:
\[
  \|x\| + \|y\| < \frac{1}{4\Ctotal}.
\]
This stricter condition ensures that the $p$-norm tail sums are bounded by a geometric
series that is controlled uniformly in $\Ctri$ (see Remark~\ref{rem:role-ctri}).

The map $(x,y)\mapsto Z(x,y)$ is continuous for the quasi-norm topology and satisfies
$e^{Z(x,y)} = e^x e^y$ in any associative quasi-Banach algebra containing $\mathfrak{g}$
as a Lie subalgebra.

If $\mathfrak{g}$ is embedded in an associative quasi-Banach algebra with
submultiplicativity constant $\Cmult$, then by Remark~\ref{rem:bracket-associative} we
have $\Cbracket \leq 2\Ctri\Cmult$ and $\Ctotal \leq 2\Ctri^2\Cmult$, yielding the
explicit sufficient condition
\[
  \|x\| + \|y\| < \frac{1}{8\Ctri^2\Cmult}.
\]
\end{theorem}

\begin{equation}\label{eq:tree}
  \|[x,y]_T\| \le (\Cbracket)^{n-1}(\|x\|+\|y\|)^n.
\end{equation}

\begin{equation}\label{eq:pnorm-bound}
  \pnorm{Z_n(x,y)}^p
  \le (2\Ctri)^p \|Z_n(x,y)\|^p
  \le (2\Ctri)^p \cdot 4^{p(n-1)}(\Cbracket)^{p(n-1)} r^{pn}.
\end{equation}

\begin{lemma}[Lipschitz continuity of BCH]\label{lem:bch-lipschitz}
Under the hypotheses of Theorem~\ref{thm:BCH}, assume further that the BCH series
converges absolutely in the sense that $\sum_{n \ge 1} \|Z_n(x,\cdot)\|$ is
term-differentiable (in particular this holds when $\mathfrak{g}$ embeds in a quasi-Banach
associative algebra and the power-series estimates of Steps~1--2 apply).
% CORRECTION: rho_0 is reduced from 1/(8*B) to 1/(16*B) so that s <= 1/(8*B),
% which is required for the Lipschitz constant L_0=2 to hold uniformly on the ball.
Set $\rho_0 := 1/(16\Cbracket)$.
For $x, y, z \in \mathfrak{g}$ with $\|x\|,\|y\|,\|z\| \le \rho_0$, the
BCH map satisfies
\[
  \|Z(x,y) - Z(x,z)\| \le \frac{1}{1 - 4\Cbracket s}\,\|y - z\|,
  \qquad s := \|x\| + \max(\|y\|,\|z\|).
\]
In particular, $Z(x,\cdot)$ is Lipschitz with constant $L_0 = 2$ on $B(0,\rho_0)$.
\end{lemma}

\begin{proof}
Write $Z(x,y) = \sum_{n\ge 1} Z_n(x,y)$ and $Z(x,z) = \sum_{n\ge 1} Z_n(x,z)$.  The
degree-$1$ component $Z_1(x,y)=x+y$ satisfies $\|Z_1(x,y)-Z_1(x,z)\| = \|y-z\|$.  For
$n\ge 2$, each $Z_n$ is a finite sum of nested commutators that are multilinear of degree
$n$ with coefficient sum $\le C_{n-1}$.

For any multilinear Lie polynomial $P(x;y)$ of total degree $n$ (degree $k$ in $x$ and
degree $n-k$ in the second argument), the inequality
\[
  \|P(x;y) - P(x;z)\| \le n \cdot (\Cbracket)^{n-1} s^{n-1} \|y-z\|,
  \quad s := \|x\| + \max(\|y\|,\|z\|),
\]
holds by multilinearity and the bracket bound \eqref{eq:tree}: each tree contributing to
$Z_n(x,\cdot)$ is multilinear, and replacing $y$ by $z$ in one leaf at a time gives a
telescoping sum of $n$ terms each bounded by $(\Cbracket)^{n-1}s^{n-1}\|y-z\|$.
Summing over the $C_{n-1}$ trees and using $n \cdot C_{n-1} \le n \cdot 4^{n-1}/n = 4^{n-1}$,
then over all $n\ge 1$ with $u:=4\Cbracket s$:
\[
  \|Z(x,y)-Z(x,z)\|
  \le \sum_{n\ge 1} n\,C_{n-1}(\Cbracket)^{n-1} s^{n-1}\|y-z\|
  \le \sum_{n\ge 1} 4^{n-1}(\Cbracket)^{n-1} s^{n-1}\|y-z\|
  = \frac{1}{1-4\Cbracket s}\|y-z\|.
\]
For $\|x\|,\|y\|,\|z\| \le \rho_0 = 1/(16\Cbracket)$, we have
$s = \|x\|+\max(\|y\|,\|z\|) \le 2\rho_0 = 1/(8\Cbracket)$,
so $4\Cbracket s \le 1/2$, giving
\[
  \|Z(x,y)-Z(x,z)\| \le \frac{1}{1 - 1/2}\|y-z\| = 2\|y-z\|,
\]
and the Lipschitz constant is $L_0 = 2$.
\end{proof}

\end{document}
